\documentclass[11pt]{article}
\usepackage{amsmath, amssymb}
\usepackage{geometry}
\usepackage{setspace}
\usepackage{booktabs}
\usepackage{enumitem}
\usepackage{fancyhdr}
\usepackage{mathtools, parskip, actuarialangle}

\geometry{margin=1in}
\setstretch{1.2}

% % Actuarial angle notation
% \newcommand{\actuarialangle}[1]{\langle #1 \rangle}

\pagestyle{fancy}
\fancyhf{}
\rhead{MATH 2790 -- Review Worksheet}
\lhead{Exam 2 Practice}
\rfoot{\thepage}

\begin{document}

\begin{center}
    {\LARGE \textbf{MATH 2790 --- Exam 2 Review Worksheet}}\\[0.5ex]
\end{center}




\vspace{1em}

%%%%%%%%%%%%%%%%%%%%%%%%%%%%%%%%%%%%%%%%%%%%%%%%%%%%%%%%%%%%%
% Problem 1 - Continuously Varying Annuity
%%%%%%%%%%%%%%%%%%%%%%%%%%%%%%%%%%%%%%%%%%%%%%%%%%%%%%%%%%%%%
% \textbf{1. Continuously Varying Annuity}\\

(1) An annuity pays continuously for 12 years at a rate of payment $R(t) = 4t + 4$ and the force of interest is 
\(\displaystyle \delta(t) = \frac{2}{t+1}\). Find the present value of this annuity.
\vspace{1em}

%%%%%%%%%%%%%%%%%%%%%%%%%%%%%%%%%%%%%%%%%%%%%%%%%%%%%%%%%%%%%
% Problem 2 - Loan Amortization
%%%%%%%%%%%%%%%%%%%%%%%%%%%%%%%%%%%%%%%%%%%%%%%%%%%%%%%%%%%%%
% \textbf{2. Loan Amortization}\\

(2) Taylor purchases a small business for \$180{,}000 financing the balance over 8 years with \emph{monthly payments} at an annual interest rate of 8.1\%.

\begin{enumerate}[label=(\alph*)]
    \item Find the level monthly payment. \textbf{Round up} your answer to 2 decimals.
    \item Compute the outstanding loan balance immediately after 50 payments. Round to 2 decimals. Notice the final payment will be adjusted so try to use the retrospective method.
\end{enumerate}

\vspace{1em}

%%%%%%%%%%%%%%%%%%%%%%%%%%%%%%%%%%%%%%%%%%%%%%%%%%%%%%%%%%%%%
% Problem 3 - Deferred Annuity Setup
%%%%%%%%%%%%%%%%%%%%%%%%%%%%%%%%%%%%%%%%%%%%%%%%%%%%%%%%%%%%%

(3) \noindent
You deposit end of month payments of $X$ into a bank account for $8$ years so that starting in the $9$th year you can withdraw $Y$ every year for $15$ years. Letting $i$ be a quarterly effective interest rate, set up an equation relating $X$ and $Y$ using actuarial notation.

\vspace{1em}

%%%%%%%%%%%%%%%%%%%%%%%%%%%%%%%%%%%%%%%%%%%%%%%%%%%%%%%%%%%%%
% Problem 4 - Increasing Annuity-Due
%%%%%%%%%%%%%%%%%%%%%%%%%%%%%%%%%%%%%%%%%%%%%%%%%%%%%%%%%%%%%
% \textbf{4. Increasing Annuity-Due}\\
(4) Ava purchases an increasing annuity-due that pays \$5 at the beginning of the first month, \$8 at the beginning of the second month, and continues to increase by \$3 per month thereafter. The last payment is \$86. The nominal annual rate is 9\% convertible monthly. Find the present value of this annuity one month before the first payment. Round to two decimals.

\vspace{1em}

%%%%%%%%%%%%%%%%%%%%%%%%%%%%%%%%%%%%%%%%%%%%%%%%%%%%%%%%%%%%%
% Problem 5 - Growing Payments (Inflation)
%%%%%%%%%%%%%%%%%%%%%%%%%%%%%%%%%%%%%%%%%%%%%%%%%%%%%%%%%%%%%
(5) You plans to deposit funds each year for 22 years to launch a start-up company right after the $22$nd deposit..  
The first deposit is \$2{,}500 and due to the excitement of the business, you increase each payment by 5\% annually.  
Assume an annual effective rate of 6.2\%. You want to throw a party exactly one year before the company launches. Find the value of \textit{all} the payments during the time of the party.


\vspace{1em}

%%%%%%%%%%%%%%%%%%%%%%%%%%%%%%%%%%%%%%%%%%%%%%%%%%%%%%%%%%%%%
% Problem 6 - Perpetuity with Rate Change
%%%%%%%%%%%%%%%%%%%%%%%%%%%%%%%%%%%%%%%%%%%%%%%%%%%%%%%%%%%%%
(6) Suppose \$40{,}000 was invested on January 1, 1980 at an annual effective interest rate of 7\% in order to provide an annual (calendar-year) scholarship of \$5{,}000 each year forever, the scholarships paid out each January 1. In what year can the first \$5{,}000 scholarship be made?
\vspace{2em}

%%%%%%%%%%%%%%%%%%%%%%%%%%%%%%%%%%%%%%%%%%%%%%%%%%%%%%%%%%%%%
% Problem 7 - Irregular Payments and Actuarial Notation
%%%%%%%%%%%%%%%%%%%%%%%%%%%%%%%%%%%%%%%%%%%%%%%%%%%%%%%%%%%%%

(7) Determine the present value of 1 payable at the end of years \(\mathbf{7, 12, 17, 22, 27, 32}\).  
Pick the correct choice and show your work to support your choice.  Leave your answer in terms of only actuarial symbols $a_{\actuarialangle{n}i}$, $s_{\actuarialangle{n}i}$, $\ddot{a}_{\actuarialangle{n}i}$, 
$\ddot{s}_{\actuarialangle{n}i}$ where $i$ is the annual effective interest rate. 
\end{document}
